\documentclass[a4paper, 12 pt]{article}
\author{}
\usepackage{color}
\usepackage[bahasa]{babel}
\usepackage{graphicx}
\title{\textbf{Rangkuman Basis Data}\linebreak}
\date{}
\begin{document}
	
	\maketitle
	\begin{center}
		\includegraphics[width=8cm,height=8cm]{D:/TEKNIK INFORMATIKA/LATEX/TUGAS DB 1/polpos.png}
	\end{center}
	\vspace{0.5 cm}
	
	\begin{center}
		Alwizain Almas Trigreisian \\
		1194004 \\
		D4TI 1A \linebreak
		\newline
		\newline
		\newline
		Untuk Memenuhi Tugas Basis Data I Semester 2 \\
		Dosen Pengampu: Syafrial Fachri Pane, ST., MTI., EBDP. \linebreak
		\newline
		\newline
		Program Studi D4 Teknik Informatika \\
		Politeknik Pos Indonesia\\
		2019/2020\\
	\end{center}
	
	\newpage
	\begin{flushleft}
		\title{\textbf{Rangkuman Basis Data}}\linebreak
	\end{flushleft}

		\par Database atau Basis Data merupakan kumpulan beberapa informasi atau data berupa gambar, angka, huruf, symbol yang saling berelasi dan berdasarkan fakta. Data-data atau informasi dalam database tersebut dibuat tabel yang memiliki attribut saling berelasi 			antara tabel satu dengan yang lainnya. Jadi konsep pada database ini seperti pada irisan himpunan di Matematika yang data-datanya saling berelasi. Adapun tempat untuk penyimpanan data atau yang disebut dengan DBMS (Database Management System) dalam 			database tersebut antara lain Oracle, Mysql, dan Firebird. Database juga memiliki tujuan, yaitu antara lain:
		\par 1. Untuk pengelompokan data
		\par 2. Mencegah data yang ganda
		\par 3. Memudahkan akses ke data
		\par 4. Menyederhanakan dari berbagai jenis dan banyaknya data\\ \linebreak

		\par Pada zaman dahulu sudah terdapat database, namun penerapannya masih manual. Contohnya seperti penyimpanan berkas data prajurit yang harus disimpan berupa lembaran-lembaran dan dikelompokkan berdasarkan jenis prajurit karena pada zaman dahulu 			belum adanya database pada sistem komputer. Namun sekarang database sudah menggunakan sistem dan sangat sederhana sehingga dapat mempermudah pengguna untuk pengelompokan data. Data-data pada database tidak diperbolehkan hanya berupa opini 			saja, namun juga harus memiliki bukti yang bisa dipertanggungjawabkan. \linebreak

		\par Dalam pembuatan program atau aplikasi yang membutuhkan database, kita harus menganalisis kebutuhan terlebih dahulu yang nantinya kebutuhan-kebutuhan tersebut	dimasukkan kedalam sistem database. Sehingga program atau aplikasi yang kita buat 	dapat 		terancang dengan baik. Banyak kasus diantara kita yang membuat sistem dahulu, baru membuat database kemudian. Hal tersebut kurang benar dalam penerapannya. Seharusnya kita membuat database terlebih dahulu lalu pembuatan sistem kemudian. Sehingga 			kebutuhan dapat dianalisa dan dapat terciptanya sistem yang baik.
\end{document}